\chapterrule
\section{规则写不出来：为什么让模型从样本中学习？}
\label{intro:00d-rules-gap}
\modelpipelinefocus{0}{手写规则与可训练模型}

我们想让程序识别猫、翻译一句话，或者回答陌生问题。
最直接的办法是把判断规则逐条写入代码。这对条件明确的任务很好用。

\subsection{先试着把规则写完}

已知直角三角形两条直角边 $a,b$，斜边满足
\[
c=\sqrt{a^2+b^2}.
\]
这条规则既明确又可检验。

但判断陌生照片是否有猫怎么办？
我们可以写“有耳朵、胡须、四条腿”，
接着遇到遮挡、背影、逆光、只拍到半张脸等情况，
又得补很多例外。问题不是猫没有规律，
而是\textbf{很难让人事先列出所有可能的图像条件}。

\subsection{把“写规则”换成“调整函数”}

另一种方法是先选择\textbf{带参数的计算函数}：
给它输入，让它产生预测，用已有样本比较预测和期望结果，
再逐步调整参数。这里学到的是适合某类数据的近似映射，
\textbf{不保证它找到了唯一的自然规律，也不保证每道新题都对}。
输入、目标、函数形式仍需要合理设计。

\subsection{从最小的数值例子开始}

暂时不引入向量、矩阵或神经网络。
我们先用只有 $w,b$ 两个参数的直线 $\widehat y=wx+b$，
让三组数字告诉我们什么叫“学习参数”。

\chaptercheck{00d-rules-gap}{为什么“猫有四条腿”不是完整的图像识别规则？}{给定三个样本 $(1,3),(2,5),(3,7)$，
试用带参数函数 $\widehat y=2x+1$ 预测 $x=4$ 时的输出；这是预测还是已知事实？}{遇到遮挡或逆光图片，应该继续硬加固定规则，
还是考虑从标注样本学习可泛化特征？说出理由。}


\begin{conclusion}
显式规则有用，但很难在某些复杂任务上列尽所有例外。
下一个缺口：如果先选带参数函数，怎样从样本中找到参数？
\end{conclusion}
